\documentclass[10pt,a4paper]{amsart}
\usepackage{amsmath,amssymb,amsthm,graphicx}
\usepackage{verbatim}
\newenvironment{reference}{\comment}{\endcomment}
%\newenvironment{reference}{}{}
\newenvironment{slogan}{\comment}{\endcomment}
\newenvironment{history}{\comment}{\endcomment}

% For commutative diagrams we use Xy-pic
\usepackage[all]{xy}

% We use 2cell for 2-commutative diagrams.
\xyoption{2cell}
% Theorem environments.
%
\theoremstyle{plain}
\newtheorem{theorem}[subsection]{Theorem}
\newtheorem{proposition}[subsection]{Proposition}
\newtheorem{lemma}[subsection]{Lemma}

\theoremstyle{definition}
\newtheorem{definition}[subsection]{Definition}
\newtheorem{example}[subsection]{Example}
\newtheorem{exercise}[subsection]{Exercise}
\newtheorem{situation}[subsection]{Situation}

\theoremstyle{remark}
\newtheorem{remark}[subsection]{Remark}
\newtheorem{remarks}[subsection]{Remarks}

\numberwithin{equation}{subsection}

% Macros
%
\def\lim{\mathop{\mathrm{lim}}\nolimits}
\def\colim{\mathop{\mathrm{colim}}\nolimits}
\def\Spec{\mathop{\mathrm{Spec}}}
\def\Hom{\mathop{\mathrm{Hom}}\nolimits}
\def\Ext{\mathop{\mathrm{Ext}}\nolimits}
\def\SheafHom{\mathop{\mathcal{H}\!\mathit{om}}\nolimits}
\def\SheafExt{\mathop{\mathcal{E}\!\mathit{xt}}\nolimits}
\def\Sch{\mathit{Sch}}
\def\Mor{\mathop{\mathrm{Mor}}\nolimits}
\def\Ob{\mathop{\mathrm{Ob}}\nolimits}
\def\Sh{\mathop{\mathit{Sh}}\nolimits}
\def\NL{\mathop{N\!L}\nolimits}
\def\CH{\mathop{\mathrm{CH}}\nolimits}
\def\proetale{{pro\text{-}\acute{e}tale}}
\def\etale{{\acute{e}tale}}
\def\QCoh{\mathit{QCoh}}
\def\Ker{\mathop{\mathrm{Ker}}}
\def\Im{\mathop{\mathrm{Im}}}
\def\Coker{\mathop{\mathrm{Coker}}}
\def\Coim{\mathop{\mathrm{Coim}}}

% Boxtimes
%
\DeclareMathSymbol{\boxtimes}{\mathbin}{AMSa}{"02}

%
% Macros for moduli stacks/spaces
%
\def\QCohstack{\mathcal{QC}\!\mathit{oh}}
\def\Cohstack{\mathcal{C}\!\mathit{oh}}
\def\Spacesstack{\mathcal{S}\!\mathit{paces}}
\def\Quotfunctor{\mathrm{Quot}}
\def\Hilbfunctor{\mathrm{Hilb}}
\def\Curvesstack{\mathcal{C}\!\mathit{urves}}
\def\Polarizedstack{\mathcal{P}\!\mathit{olarized}}
\def\Complexesstack{\mathcal{C}\!\mathit{omplexes}}
% \Pic is the operator that assigns to X its picard group, usage \Pic(X)
% \Picardstack_{X/B} denotes the Picard stack of X over B
% \Picardfunctor_{X/B} denotes the Picard functor of X over B
\def\Pic{\mathop{\mathrm{Pic}}\nolimits}
\def\Picardstack{\mathcal{P}\!\mathit{ic}}
\def\Picardfunctor{\mathrm{Pic}}
\def\Deformationcategory{\mathcal{D}\!\mathit{ef}}
\usepackage{cleveref}
\crefname{originalreference}{original reference}{original references}
\Crefname{originalreference}{Original reference}{Original references}
\makeatletter
\newlabel{lemma-galois-profinite}{{0BMJ}{source}{}{sourcepointer-1}{}}
\newlabel{lemma-galois-profinite@cref}{{[originalreference][0][]0BMJ}{[1][0][]source}{}{}{}}
\makeatother
\crefname{originalreference}{original reference}{original references}
\Crefname{originalreference}{Original reference}{Original references}
\makeatletter
\newlabel{lemma-lift-maps}{{0BME}{source}{}{sourcepointer-2}{}}
\newlabel{lemma-lift-maps@cref}{{[originalreference][0][]0BME}{[1][0][]source}{}{}{}}
\makeatother
\title{Restriction gives a canonical continuous surjection of Galois groups in a tower}
\author{The Stacks Project Authors}
\date{stacks-project:fields.tex:local-snapshot-20261012; GNU FDL 1.2 or later}
\begin{document}
\maketitle
\noindent\textit{Editorial note.} All delivered mathematical text is the author's own, copied from the pinned Stacks Project chapter file fields.tex; no statement or proof step was rewritten, reordered or removed.\par\medskip
\noindent\textit{Editorial note.} The counted proof is the author's complete argument: The restriction map of this lemma is exactly the map that the exact-sequence restatement (declined unit lemma-ses-infinite-galois) would have packaged.\par\medskip
\noindent\textit{Editorial note.} The label lemma-galois-infinite names the counted statement, and the proof environment immediately following it in the source is delivered verbatim as its proof.\par\medskip
\noindent\textit{Editorial note.} The counted statement is carried by the author's own lemma environment in the chapter source.\par\medskip
\noindent\textit{Allowed standard prerequisites (editorial).}\begin{itemize}
\item Chapter lemma-lift-maps (delivered card L48358): item 1 restricts automorphisms to normal subextensions; item 2 extends K-maps into a normal extension to automorphisms
\item Chapter lemma-galois-profinite (delivered card L48328): part (1) is the universal property of the profinite topology, characterizing continuity of group maps through continuity of the action
\item Chapter section setup: Gal(L over K) carries the Krull topology, for which the action Gal(L over K) x L to L is continuous
\end{itemize}\medskip
\ifdefined\hypertarget\hypertarget{sourcepointer-1}{}\fi
\noindent\textit{Source reference.} Original reference 0BMJ: sources/stacks-fields/pinned-original-source.tex, source lines 3076--3076. Prerequisite pointer invoked at 3187; it names a standard fact of the same chapter that the counted argument uses, and every step of the counted proof built on it is retained verbatim in the delivered body.\par\medskip
\ifdefined\hypertarget\hypertarget{sourcepointer-2}{}\fi
\noindent\textit{Source reference.} Original reference 0BME: sources/stacks-fields/pinned-original-source.tex, source lines 2163--2163. Prerequisite pointer invoked at 3175; it names a standard fact of the same chapter that the counted argument uses, and every step of the counted proof built on it is retained verbatim in the delivered body.\par\medskip
\begin{lemma}
\label{lemma-galois-infinite}
Let $L/M/K$ be a tower of fields. Assume both $L/K$ and
$M/K$ are Galois. Then there is a canonical surjective continuous
homomorphism $c : \text{Gal}(L/K) \to \text{Gal}(M/K)$.
\end{lemma}

\begin{proof}
By Lemma \ref{lemma-lift-maps} given $\tau : L \to L$ in
$\text{Gal}(L/K)$ the restriction $\tau|_M : M \to M$
is an element of $\text{Gal}(M/K)$. This defines the homomorphism $c$.
Continuity follows from the universal property of the topology:
the action
$$
\text{Gal}(L/K) \times M \longrightarrow M,\quad
(\tau, x) \longmapsto \tau(x) = c(\tau)(x)
$$
is continuous as $M \subset L$ and the action
$\text{Gal}(L/K) \times L \to L$ is continuous.
Hence continuity of $c$ by part (1) of
Lemma \ref{lemma-galois-profinite}.
Lemma \ref{lemma-lift-maps} also
shows that the map is surjective.
\end{proof}

\end{document}
